\documentclass[11pt]{article}

\usepackage{amsmath,amssymb,amsthm,mathtools}
\usepackage[margin=1in]{geometry}
\usepackage{microtype}
\usepackage[hidelinks]{hyperref}

\newtheorem{theorem}{Theorem}[section]
\newtheorem{lemma}[theorem]{Lemma}
\newtheorem{proposition}[theorem]{Proposition}
\newtheorem{corollary}[theorem]{Corollary}
\newtheorem{remark}[theorem]{Remark}

\newcommand{\dd}{\,\mathrm{d}}
\newcommand{\one}{\mathbf{1}}

\title{The Sidorenko Inequality for Forests}
\author{}
\date{}

\begin{document}

\maketitle

\begin{abstract}
Let $F$ be a finite forest with $e\geq 1$ edges, and let $V$ be a
$[0,1]$-valued symmetric measurable kernel on an arbitrary probability
space.  We prove from first principles that
\[
t(F,V)\geq t(K_2,V)^e.
\]
The connected case is proved by a leaf-removal induction which constructs
a probability density on the assignments of a tree.  Its one-vertex
marginals are the degree-biased measure.  Two explicit applications of
Jensen's inequality then give the desired bound.  Disjoint components,
isolated vertices, and zero edge density are treated separately.
\end{abstract}


%%%%%%%%%%%%%%%%%%%%%%%%%%%%%%%%%%%%%%%%%%%%%%%%%%%%%%%%%%%%%%%%%%%%%%%%%%%%%%%
\section{Definitions and elementary reductions}
%%%%%%%%%%%%%%%%%%%%%%%%%%%%%%%%%%%%%%%%%%%%%%%%%%%%%%%%%%%%%%%%%%%%%%%%%%%%%%%

Let $(\Omega,\mu)$ be a probability space, with no further assumptions
on its measurable structure.  Let
\[
V\colon \Omega^2\longrightarrow[0,1]
\]
be measurable with respect to the product sigma-algebra and symmetric.
For a finite simple graph $G$, define its ordinary, non-injective
homomorphism density by
\begin{equation}
t(G,V)
=
\int_{\Omega^{V(G)}}
\prod_{uv\in E(G)}V(x_u,x_v)
\prod_{u\in V(G)}\dd\mu(x_u).
\label{eq:hom-density}
\end{equation}
The integrand is measurable, nonnegative, and at most $1$.  Consequently
all integrals in \eqref{eq:hom-density} exist and are finite.  Every
reordering or partial evaluation of these finite product integrals below
is justified by the Fubini--Tonelli theorem; its applicability follows
either from nonnegativity or, more strongly, from boundedness by $1$.

Set
\begin{equation}
d(x):=\int_\Omega V(x,y)\dd\mu(y),
\qquad
z:=t(K_2,V)=\int_\Omega d(x)\dd\mu(x).
\label{eq:degree}
\end{equation}
Tonelli's theorem shows that $d$ is measurable and gives the second
identity.  Since $0\leq V\leq1$ and $\mu(\Omega)=1$, one has
$0\leq d\leq1$ and $0\leq z\leq1$.

We first record precisely how vertices and components that do not
interact in the graph behave in the homomorphism integral.

\begin{lemma}[Isolated vertices and disjoint unions]
\label{lem:components}
Let $G_1$ and $G_2$ be finite simple graphs with disjoint vertex sets.
Then
\begin{equation}
t(G_1\sqcup G_2,V)=t(G_1,V)t(G_2,V).
\label{eq:multiplicativity}
\end{equation}
In particular, adjoining or deleting an isolated vertex does not change
homomorphism density.
\end{lemma}

\begin{proof}
Write
\[
f_i((x_u)_{u\in V(G_i)})
:=\prod_{uv\in E(G_i)}V(x_u,x_v)
\qquad (i=1,2).
\]
The integrand for $G_1\sqcup G_2$ is $f_1(x)f_2(y)$.  The finite product
measure on its vertex set is the product of the finite product measures
on $V(G_1)$ and $V(G_2)$.  Fubini--Tonelli applies because
$f_1f_2$ is measurable, nonnegative, and bounded by $1$, and gives
\begin{align*}
t(G_1\sqcup G_2,V)
&=\int_{\Omega^{V(G_1)}}\int_{\Omega^{V(G_2)}}
f_1(x)f_2(y)\prod_{v\in V(G_2)}\dd\mu(y_v)
\prod_{u\in V(G_1)}\dd\mu(x_u)
\\
&=\int_{\Omega^{V(G_1)}}
f_1(x)t(G_2,V)\prod_{u\in V(G_1)}\dd\mu(x_u)
\\
&=t(G_1,V)t(G_2,V),
\end{align*}
which is \eqref{eq:multiplicativity}.  For the one-vertex graph $K_1$,
the edge product is the empty product $1$, so
\[
t(K_1,V)=\int_\Omega1\dd\mu=\mu(\Omega)=1.
\]
Applying \eqref{eq:multiplicativity} with $G_2=K_1$ proves the last
assertion.
\end{proof}

The case of zero edge density will not require division by $z$.

\begin{lemma}[Zero edge density]
\label{lem:zero}
If $z=0$ and $G$ has at least one edge, then $t(G,V)=0$.
\end{lemma}

\begin{proof}
Fix an edge $ab\in E(G)$.  Every factor in the integrand of
\eqref{eq:hom-density} lies in $[0,1]$, so pointwise
\[
0\leq \prod_{uv\in E(G)}V(x_u,x_v)\leq V(x_a,x_b).
\]
Integrating this inequality over $\Omega^{V(G)}$ is legitimate by
monotonicity of the integral.  Fubini--Tonelli applies to the right-hand
side because it is nonnegative and bounded, and the integrations in all
coordinates other than $a,b$ contribute the factor
$\mu(\Omega)^{|V(G)|-2}=1$.  Hence
\[
0\leq t(G,V)\leq
\int_{\Omega^2}V(x_a,x_b)\dd\mu(x_a)\dd\mu(x_b)=z=0.
\]
Thus $t(G,V)=0$.
\end{proof}


%%%%%%%%%%%%%%%%%%%%%%%%%%%%%%%%%%%%%%%%%%%%%%%%%%%%%%%%%%%%%%%%%%%%%%%%%%%%%%%
\section{A probability density attached to a tree}
%%%%%%%%%%%%%%%%%%%%%%%%%%%%%%%%%%%%%%%%%%%%%%%%%%%%%%%%%%%%%%%%%%%%%%%%%%%%%%%

For this section assume $z>0$.  Define a probability measure $\nu$ on
$\Omega$ by
\begin{equation}
\dd\nu(x)=\frac{d(x)}{z}\dd\mu(x).
\label{eq:nu}
\end{equation}
Indeed, $d/z$ is nonnegative and measurable, and
\eqref{eq:degree} gives $\int(d/z)\dd\mu=1$.  Notice also that
\begin{equation}
\nu(\{x:d(x)=0\})
=\frac1z\int_{\{d=0\}}d(x)\dd\mu(x)=0.
\label{eq:positive-degree-nu}
\end{equation}

Let $T$ be a finite tree with $m\geq1$ edges.  Every vertex of $T$ has
positive degree.  Define a nonnegative measurable function
$g_T$ on $\Omega^{V(T)}$ by
\begin{equation}
g_T((x_v)_{v\in V(T)})
:=
\frac{\displaystyle\prod_{uv\in E(T)}V(x_u,x_v)}
{\displaystyle z\prod_{v\in V(T)}d(x_v)^{\deg_T(v)-1}}
\label{eq:tree-density}
\end{equation}
when every factor in the denominator having a positive exponent is
positive, and set $g_T=0$ otherwise.  Exponents equal to zero contribute
the factor $1$, including when $d(x_v)=0$.

The next proposition is the inductive part of the proof.  It establishes
both the normalization of $g_T$ and every marginal that will be needed
later.

\begin{proposition}[Tree-walk normalization and marginals]
\label{prop:tree-density}
For every finite tree $T$ with at least one edge,
\begin{equation}
\int_{\Omega^{V(T)}}g_T\prod_{v\in V(T)}\dd\mu(x_v)=1.
\label{eq:normalization}
\end{equation}
Moreover, for every vertex $a\in V(T)$ and every measurable set
$B\subseteq\Omega$,
\begin{equation}
\int_{\Omega^{V(T)}}
\one_B(x_a)g_T(x)
\prod_{v\in V(T)}\dd\mu(x_v)
=\nu(B).
\label{eq:marginal}
\end{equation}
Thus the probability measure with density $g_T$ relative to
$\mu^{V(T)}$ has one-vertex marginal $\nu$ at every vertex.
\end{proposition}

\begin{proof}
We use induction on the positive integer
\[
m=|E(T)|.
\]
The precise induction assertion $P(m)$ is: for every finite tree with
exactly $m$ edges, the function in \eqref{eq:tree-density} satisfies
\eqref{eq:normalization} and \eqref{eq:marginal} for every one of its
vertices and every measurable $B\subseteq\Omega$.  The measure is the
number of edges; in the induction step it decreases from $m$ to $m-1$,
so the induction is well-founded on the positive integers.

For the base case $m=1$, let the endpoints be $a,b$.  Both degrees are
$1$, and hence
\[
g_T(x_a,x_b)=\frac{V(x_a,x_b)}{z}.
\]
Its integral is $1$ by the definition of $z$.  For a measurable
$B\subseteq\Omega$, Fubini--Tonelli applies because the following
integrands are nonnegative, and gives
\begin{align*}
\int_{\Omega^2}\one_B(x_a)g_T(x_a,x_b)
\dd\mu(x_a)\dd\mu(x_b)
&=\frac1z\int_B d(x_a)\dd\mu(x_a)
=\nu(B).
\end{align*}
The same computation for $b$ uses symmetry of $V$:
\begin{align*}
\int_{\Omega^2}\one_B(x_b)g_T(x_a,x_b)
\dd\mu(x_a)\dd\mu(x_b)
&=\frac1z\int_B
\left(\int_\Omega V(x_a,x_b)\dd\mu(x_a)\right)\dd\mu(x_b)
\\
&=\frac1z\int_B d(x_b)\dd\mu(x_b)
=\nu(B).
\end{align*}
This proves $P(1)$.

Now let $m\geq2$ and assume $P(m-1)$.  Choose a longest simple path in
$T$, let $\ell$ be one of its endpoints, and let $p$ be the vertex next
to $\ell$ on that path.  The endpoint $\ell$ has no neighbour outside
the path: such a neighbour would extend the path without repeating a
vertex, contrary to maximality.  Thus $\ell$ is a leaf and $p$ is its
unique neighbour.  Delete $\ell$ and its incident edge $p\ell$, and
call the resulting graph $T'$.  Deleting a leaf preserves connectedness
among the remaining vertices and cannot create a cycle.  Therefore
$T'$ is a tree with exactly $m-1\geq1$ edges, so $P(m-1)$ applies to it.

Define
\begin{equation}
R(x,y):=
\begin{cases}
V(x,y)/d(x),&d(x)>0,\\
0,&d(x)=0.
\end{cases}
\label{eq:transition-density}
\end{equation}
This is a nonnegative measurable function.  Direct substitution of
the degrees
\[
\deg_T(p)=\deg_{T'}(p)+1,
\qquad
\deg_T(v)=\deg_{T'}(v)\quad(v\in V(T')\setminus\{p\}),
\qquad
\deg_T(\ell)=1
\]
into \eqref{eq:tree-density} gives the leaf factorization
\begin{equation}
g_T(x)=g_{T'}(x|_{V(T')})R(x_p,x_\ell).
\label{eq:leaf-factorization}
\end{equation}
If $d(x_p)>0$, this is cancellation of exactly one additional factor
$d(x_p)$ in the denominator.  If $d(x_p)=0$, both sides are zero by
the definitions of $g_T$ and $R$.  Thus \eqref{eq:leaf-factorization}
holds pointwise.

For every $x\in\Omega$, the definition of $d$ gives
\begin{equation}
\int_\Omega R(x,y)\dd\mu(y)=\one_{\{d>0\}}(x).
\label{eq:R-mass}
\end{equation}
Using \eqref{eq:leaf-factorization}, Tonelli's theorem for the
nonnegative integrand, \eqref{eq:R-mass}, and then the induction
hypothesis for the $p$-marginal, we obtain
\begin{align*}
\int_{\Omega^{V(T)}}g_T\dd\mu^{V(T)}
&=\int_{\Omega^{V(T')}}g_{T'}(x)
\left(\int_\Omega R(x_p,y)\dd\mu(y)\right)
\dd\mu^{V(T')}(x)
\\
&=\int_{\Omega^{V(T')}}g_{T'}(x)
\one_{\{d>0\}}(x_p)\dd\mu^{V(T')}(x)
\\
&=\nu(\{d>0\})=1,
\end{align*}
where the last equality is \eqref{eq:positive-degree-nu}.  This proves
normalization.

Let $a\in V(T')$ and let $B\subseteq\Omega$ be measurable.  The same
calculation, with the factor $\one_B(x_a)$ retained, gives
\begin{align}
\int_{\Omega^{V(T)}}\one_B(x_a)g_T\dd\mu^{V(T)}
&=\int_{\Omega^{V(T')}}\one_B(x_a)g_{T'}(x)
\one_{\{d>0\}}(x_p)\dd\mu^{V(T')}(x).
\label{eq:old-marginal-first}
\end{align}
The part of the last integral on $\{d(x_p)=0\}$ is zero: by the
induction hypothesis, its value without the factor
$\one_B(x_a)\leq1$ is the $p$-marginal mass
$\nu(\{d=0\})=0$.  Removing $\one_{\{d>0\}}(x_p)$ from
\eqref{eq:old-marginal-first} and applying the induction hypothesis to
the $a$-marginal therefore gives
\[
\int_{\Omega^{V(T)}}\one_B(x_a)g_T\dd\mu^{V(T)}=\nu(B).
\]

It remains to calculate the new leaf marginal.  Define
\[
\psi_B(x):=\int_B R(x,y)\dd\mu(y).
\]
Tonelli's theorem applies at each step below because all integrands are
nonnegative.  The factorization \eqref{eq:leaf-factorization}, followed
by the induction hypothesis for the $p$-marginal, yields
\begin{align}
\int_{\Omega^{V(T)}}\one_B(x_\ell)g_T\dd\mu^{V(T)}
&=\int_{\Omega^{V(T')}}g_{T'}(x)\psi_B(x_p)
\dd\mu^{V(T')}(x)
\\
&=\int_\Omega\psi_B(x)\dd\nu(x)
\\
&=\frac1z\int_{\{d>0\}}\int_B V(x,y)\dd\mu(y)\dd\mu(x).
\label{eq:leaf-marginal-intermediate}
\end{align}
The omitted part with $d(x)=0$ has integral zero, because
\[
0\leq\int_{\{d=0\}}\int_B V(x,y)\dd\mu(y)\dd\mu(x)
\leq\int_{\{d=0\}}d(x)\dd\mu(x)=0.
\]
After adding that zero part to \eqref{eq:leaf-marginal-intermediate},
symmetry of $V$ and Fubini--Tonelli give
\begin{align*}
\int_{\Omega^{V(T)}}\one_B(x_\ell)g_T\dd\mu^{V(T)}
&=\frac1z\int_B\int_\Omega V(x,y)\dd\mu(x)\dd\mu(y)
\\
&=\frac1z\int_B d(y)\dd\mu(y)
=\nu(B).
\end{align*}
Thus every vertex marginal is $\nu$, completing the proof of $P(m)$
and the induction.
\end{proof}


%%%%%%%%%%%%%%%%%%%%%%%%%%%%%%%%%%%%%%%%%%%%%%%%%%%%%%%%%%%%%%%%%%%%%%%%%%%%%%%
\section{The inequality for a tree}
%%%%%%%%%%%%%%%%%%%%%%%%%%%%%%%%%%%%%%%%%%%%%%%%%%%%%%%%%%%%%%%%%%%%%%%%%%%%%%%

We now compare the tree-walk probability measure with the original
homomorphism integral.  The logarithms used below are harmless even when
$d$ approaches zero; the required integrability is verified before
Jensen's inequality is applied.

\begin{theorem}[Tree case]
\label{thm:tree}
If $T$ is a finite tree with $m\geq1$ edges, then
\[
t(T,V)\geq z^m.
\]
\end{theorem}

\begin{proof}
The conclusion for $z=0$ is Lemma~\ref{lem:zero}, so suppose $z>0$.
Let $Q_T$ be the probability measure on $\Omega^{V(T)}$ whose density
with respect to $\mu^{V(T)}$ is $g_T$; it is a probability measure by
Proposition~\ref{prop:tree-density}.

Define
\begin{equation}
A_T(x):=z\prod_{v\in V(T)}d(x_v)^{\deg_T(v)-1}.
\label{eq:A}
\end{equation}
We claim that
\begin{equation}
t(T,V)=\int_{\Omega^{V(T)}}A_T(x)\dd Q_T(x).
\label{eq:t-as-expectation}
\end{equation}
On the set on which all positive-exponent factors in \eqref{eq:A} are
positive, the identity
\[
A_T(x)g_T(x)=\prod_{uv\in E(T)}V(x_u,x_v)
\]
is immediate from \eqref{eq:tree-density}.  On the complementary set,
choose a vertex $a$ such that $\deg_T(a)-1>0$ and $d(x_a)=0$.  The
integral of the homomorphism integrand over this set is zero: choose an
edge $ab$ incident to $a$ and use
\begin{align*}
0
&\leq
\int_{\Omega^{V(T)}}
\one_{\{d(x_a)=0\}}
\prod_{uv\in E(T)}V(x_u,x_v)\dd\mu^{V(T)}(x)
\\
&\leq
\int_{\Omega^{V(T)}}
\one_{\{d(x_a)=0\}}V(x_a,x_b)\dd\mu^{V(T)}(x)
\\
&=\int_{\{d=0\}}d(x_a)\dd\mu(x_a)=0.
\end{align*}
Here the inequality uses $0\leq V\leq1$, and the final evaluation uses
Fubini--Tonelli for a bounded nonnegative integrand.  There are only
finitely many possible vertices $a$, so the union of these exceptional
sets also contributes zero.  On that union, $A_Tg_T=0$ by definition.
This proves \eqref{eq:t-as-expectation}.

For $s\in[0,1]$, set $0\log0:=0$.  The elementary bound
$0\leq-s\log s\leq1/\mathrm{e}$ shows that $d\log d$ is integrable with
respect to $\mu$.  By the marginal assertion in
Proposition~\ref{prop:tree-density},
\begin{equation}
\int_\Omega |\log d(x)|\dd\nu(x)
=\frac1z\int_\Omega -d(x)\log d(x)\dd\mu(x)
<\infty,
\label{eq:log-integrable}
\end{equation}
where the integrand is assigned value $0$ at $d=0$; this assignment is
valid because $\nu(\{d=0\})=0$.  Since every exponent in
\eqref{eq:A} is nonnegative, \eqref{eq:log-integrable} and the
one-vertex marginals show that $\log A_T$ is integrable with respect to
$Q_T$.  Moreover, $0<A_T\leq z\leq1$ holds $Q_T$-almost everywhere:
positive-exponent zero factors form a $Q_T$-null set by the marginals,
and all degree values are at most $1$.  In particular,
$\exp(\log A_T)=A_T$ is $Q_T$-integrable.

Apply Jensen's inequality to the convex function
$r\mapsto\exp(r)$ with respect to the probability measure $Q_T$.
The preceding paragraph verifies that both $\log A_T$ and its
exponential are integrable.  Using \eqref{eq:t-as-expectation}, we get
\begin{equation}
t(T,V)=\int A_T\dd Q_T
=\int\exp(\log A_T)\dd Q_T
\geq\exp\left(\int\log A_T\dd Q_T\right).
\label{eq:exp-jensen}
\end{equation}
Every one-vertex marginal of $Q_T$ is $\nu$, so
\begin{align}
\int\log A_T\dd Q_T
&=\log z+
\sum_{v\in V(T)}(\deg_T(v)-1)
\int_\Omega\log d(x)\dd\nu(x)
\notag\\
&=\log z+
\left(\sum_{v\in V(T)}(\deg_T(v)-1)\right)
\frac1z\int_\Omega d(x)\log d(x)\dd\mu(x).
\label{eq:logA}
\end{align}
All terms are finite by \eqref{eq:log-integrable}, which also justifies
the finite sum and the marginal substitutions.

We next apply Jensen's inequality to the function
\[
\phi(s)=
\begin{cases}
s\log s,&s>0,\\
0,&s=0.
\end{cases}
\]
This function is continuous and convex on $[0,1]$: on $(0,1]$ one has
$\phi''(s)=1/s>0$, and continuity at $0$ extends convexity to the closed
interval.  The measure in this application is the original probability
measure $\mu$.  The random variable $d$ takes values in $[0,1]$, and
$\phi(d)$ is integrable because $|d\log d|\leq1/\mathrm{e}$.  Hence
Jensen gives
\begin{equation}
\int_\Omega d(x)\log d(x)\dd\mu(x)
\geq
\left(\int_\Omega d(x)\dd\mu(x)\right)
\log\left(\int_\Omega d(x)\dd\mu(x)\right)
=z\log z.
\label{eq:xlogx-jensen}
\end{equation}

A finite tree with $m$ edges has $m+1$ vertices, and the handshake
identity gives $\sum_v\deg_T(v)=2m$.  Therefore
\begin{equation}
\sum_{v\in V(T)}(\deg_T(v)-1)
=2m-(m+1)=m-1.
\label{eq:degree-sum}
\end{equation}
The coefficient $m-1$ is nonnegative.  Dividing
\eqref{eq:xlogx-jensen} by $z>0$, substituting it and
\eqref{eq:degree-sum} into \eqref{eq:logA}, and then using the fact that
the exponential function is increasing, we obtain
\begin{align*}
\int\log A_T\dd Q_T
&\geq\log z+(m-1)\log z=m\log z,
\\
\exp\left(\int\log A_T\dd Q_T\right)
&\geq\exp(m\log z)=z^m.
\end{align*}
Combining this with \eqref{eq:exp-jensen} proves the theorem.
\end{proof}


%%%%%%%%%%%%%%%%%%%%%%%%%%%%%%%%%%%%%%%%%%%%%%%%%%%%%%%%%%%%%%%%%%%%%%%%%%%%%%%
\section{Forests}
%%%%%%%%%%%%%%%%%%%%%%%%%%%%%%%%%%%%%%%%%%%%%%%%%%%%%%%%%%%%%%%%%%%%%%%%%%%%%%%

\begin{theorem}[Sidorenko inequality for forests]
\label{thm:forest}
Let $F$ be a finite forest with $e=e(F)\geq1$ edges.  On every
probability space $(\Omega,\mu)$, every symmetric measurable kernel
$V\colon\Omega^2\to[0,1]$ satisfies
\[
\boxed{t(F,V)\geq z^e,\qquad z=t(K_2,V).}
\]
\end{theorem}

\begin{proof}
If $z=0$, the forest $F$ has at least one edge by assumption, so
Lemma~\ref{lem:zero} gives $t(F,V)=0=z^e$.

Suppose now that $z>0$.  Decompose $F$ into its finitely many connected
components.  Let $T_1,\ldots,T_k$ be exactly the components containing
an edge, and let $r$ be the number of isolated-vertex components.  The
assumption $e\geq1$ implies $k\geq1$.  Since $F$ is acyclic, every
connected component containing an edge is a finite tree.  Write
$e_i=|E(T_i)|\geq1$.  Every edge belongs to exactly one component, so
\begin{equation}
e=\sum_{i=1}^k e_i.
\label{eq:edge-decomposition}
\end{equation}

Repeated application of Lemma~\ref{lem:components}, including its
isolated-vertex assertion, gives
\begin{equation}
t(F,V)=\left(\prod_{i=1}^k t(T_i,V)\right)t(K_1,V)^r
=\prod_{i=1}^k t(T_i,V).
\label{eq:forest-factorization}
\end{equation}
Theorem~\ref{thm:tree} gives $t(T_i,V)\geq z^{e_i}$ for each $i$.
All quantities are nonnegative, so multiplying these finitely many
inequalities and then using \eqref{eq:edge-decomposition} yields
\[
t(F,V)
=\prod_{i=1}^k t(T_i,V)
\geq\prod_{i=1}^k z^{e_i}
=z^{\sum_{i=1}^k e_i}
=z^e.
\]
This proves the stated inequality for disconnected forests and shows
explicitly that isolated vertices affect neither side: they contribute
a multiplicative factor $1$ to the left and no edges to the exponent on
the right.
\end{proof}


%%%%%%%%%%%%%%%%%%%%%%%%%%%%%%%%%%%%%%%%%%%%%%%%%%%%%%%%%%%%%%%%%%%%%%%%%%%%%%%
\section{Scope and external inputs}
%%%%%%%%%%%%%%%%%%%%%%%%%%%%%%%%%%%%%%%%%%%%%%%%%%%%%%%%%%%%%%%%%%%%%%%%%%%%%%%

The proof is self-contained apart from the following facts, listed in
the precise forms in which they are used.

\small

\begin{itemize}
\item \textbf{Finite product measures and Fubini--Tonelli.}
Finite products of a probability measure exist on the corresponding
product sigma-algebras, may be regrouped according to a partition of
the finite coordinate set, and are again probability measures.
For a nonnegative measurable function on such a finite product, iterated
integrals may be taken in any order and equal the product-space integral;
the partial integral is measurable.  The same conclusion holds for an
integrable function.  We use the nonnegative form for homomorphism
integrands, the densities $g_T$, their indicator-weighted versions, and
the leaf factor $R$.  All homomorphism integrands are also bounded by
$1$, so their integrability is automatic.

\item \textbf{Elementary integration facts.}
The integral is monotone on nonnegative measurable functions; a
nonnegative measurable function with integral zero vanishes almost
everywhere; finite unions of null sets are null; and functions equal
almost everywhere have equal integrals when either is nonnegative or
integrable.  A measurable density $g$ of integral $1$ defines a
probability measure $Q$ by $\dd Q=g\dd P$, with
$\int h\dd Q=\int hg\dd P$.  Equality on all measurable-set indicators
identifies two measures, which then integrate every nonnegative
measurable or integrable function equally.  These facts are used in the
zero-density argument, on degree-zero sets, and for $g_T$ and $Q_T$.

\item \textbf{Jensen's inequality.}
If $P$ is a probability measure, $X$ and $\varphi(X)$ are integrable,
and $\varphi$ is convex on an interval containing the range of $X$, then
\[
\varphi\!\left(\int X\dd P\right)
\leq\int\varphi(X)\dd P.
\]
It is used first with $P=Q_T$, $X=\log A_T$, and
$\varphi(r)=e^r$, after integrability is verified in
\eqref{eq:log-integrable}; it is used second with $P=\mu$, $X=d$, and
$\varphi(s)=s\log s$ with $\varphi(0)=0$, after boundedness and
integrability are verified immediately before
\eqref{eq:xlogx-jensen}.  No application of Cauchy--Schwarz or
H\"older's inequality is used.

\item \textbf{Elementary calculus and real arithmetic.}
The exponential function is convex and increasing; the function
$s\mapsto s\log s$, extended by $0$ at $s=0$, is continuous and convex
on $[0,1]$; $0\leq-s\log s\leq1/\mathrm{e}$ on $[0,1]$; and the usual
finite-sum, power, logarithm, and exponential identities hold for
positive real numbers.  These are used only in the two Jensen steps and
the displayed algebra following them.  We also use that finite products
of nonnegative inequalities may be multiplied without reversing the
inequality, and that a product of numbers in $[0,1]$ lies in $[0,1]$
and is at most any selected factor.

\item \textbf{Elementary finite graph facts.}
Every finite graph is the disjoint union of its connected components;
every edge lies in exactly one component; a connected acyclic finite
graph with an edge is a tree; a finite tree with $m$ edges has $m+1$
vertices; and the handshake identity is
$\sum_v\deg(v)=2m$.  A longest simple path in a finite tree has leaf
endpoints, and deleting one endpoint and its incident edge leaves a tree
with one fewer edge.  These facts are used respectively in the forest
reduction, \eqref{eq:degree-sum}, and the induction in
Proposition~\ref{prop:tree-density}.  We also use the well-founded
induction principle on the positive integers.

\item \textbf{Measurability closure.}
Coordinate projections, finite products and quotients on measurable
positive-denominator sets of real-valued measurable functions are
measurable.  Together with Fubini--Tonelli, this ensures that the degree
function $d$, the homomorphism integrands, $R$, $g_T$, and $A_T$ are
measurable.
\end{itemize}

For the requested scope check: (a) no step assumes that $\Omega$ is
atomless, standard Borel, separable, or countably generated; only its
probability measure and finite product measures are used.  (b) No step
assumes that $V$ or $d$ is bounded away from zero; degree-zero sets are
handled explicitly in \eqref{eq:positive-degree-nu},
\eqref{eq:transition-density}, and \eqref{eq:t-as-expectation}.
(c) The induction is on the positive integer $|E(T)|$, decreases by one
after leaf removal, terminates at $|E(T)|=1$, and that base case is
computed explicitly.  (d) Disconnected forests and isolated vertices
are covered by Lemma~\ref{lem:components} and
\eqref{eq:forest-factorization}.  (e) The case $z=0$ is proved in
Lemma~\ref{lem:zero} and invoked at the start of both
Theorem~\ref{thm:tree} and Theorem~\ref{thm:forest}.  No finite-graph
limit, regularity lemma, weak-regularity argument, compactness argument,
or graphon limit theory is used.

\normalsize

%%%%%%%%%%%%%%%%%%%%%%%%%%%%%%%%%%%%%%%%%%%%%%%%%%%%%%%%%%%%%%%%%%%%%%%%%%%%%%%
\appendix
\section{What the Lean formalization actually proves}
%%%%%%%%%%%%%%%%%%%%%%%%%%%%%%%%%%%%%%%%%%%%%%%%%%%%%%%%%%%%%%%%%%%%%%%%%%%%%%%

Neither Theorem~\ref{thm:tree} nor Theorem~\ref{thm:forest} is formalized.  The
leaf-removal induction, the density \eqref{eq:tree-density} on tree
assignments, and its marginal computation have no counterpart in the Lean
development.

What the catalogue needs from forest Sidorenko is only the instances that the
cone-over-a-forest rows consume, and each is available by an elementary route
that avoids the induction:

\begin{itemize}
\item stars $K_{1,n}$, where $t(K_{1,n},V)=M_n\geq z^n$ is Jensen for
 $t\mapsto t^n$ against the degree function;
\item disjoint unions of single edges, by multiplicativity, with equality;
\item $P_4$, and only $P_4$, which is
 \texttt{Methods/PathSidorenko.lean}: three-factor H\"older on the
 almost-everywhere factorization
 $W=(Wdd')^{1/3}(W/d)^{1/3}(W/d')^{1/3}$, the second and third factors
 integrating to $\mu\{d>0\}\leq1$.
\end{itemize}

So the general theorem is stated nowhere in Lean, and a row needing a tree
other than a star or $P_4$ would need it proved first.

\end{document}
